\item[\underline{\underline{Transaction-row n$^°~(\bm{i + \roffTxSkipUserTransaction})$:}}]
	we define shorthands
	\[
		\left\{ \begin{array}{lcl}
			\locIsDeploymentTransaction  & \define & \txIsDeployment _{i + \roffTxSkipUserTransaction} \\
			\locIsMessageCallTransaction & \define & 1 - \locIsDeploymentTransaction                   \\
		\end{array} \right.
	\]
	we peek into the transaction and justify portions of it:
	\begin{description}
		\item[\underline{(Partially) justifying \txRequiresEvmExecution{}:}]
			we impose the following
			\begin{enumerate}
				\item $\txRequiresEvmExecution_{i + \roffTxSkipUserTransaction} = \rZero$;
				\item
					\If   $\locIsMessageCallTransaction = 1$
					\Then
					\[
						\left[ \begin{array}{cl}
							+ & \locRecipientHasEmptyCode                       \\
							+ & \locRecipientIsDelegatedAndDelegateHasEmptyCode \\
							% + & \locRecipientIsDelegatedAndDelegateIsDelegated  \\
						\end{array} \right]
						= 1
					\]
					where we set
					\[
						\locRecipientIsDelegatedAndDelegateHasEmptyCode \define
						\left[ \begin{array}{clr}
							\cdot & \locRecipientIsDelegated \\
							\cdot & \locDelegateHasEmptyCode \\
						\end{array} \right]
						% \hspace*{-2cm}
						% \left\{ \begin{array}{lcl}
						% 	\locRecipientIsDelegatedAndDelegateHasEmptyCode & \define &
						% 	\left[ \begin{array}{clr}
						% 		\cdot & \locRecipientIsDelegated \\
						% 		\cdot & \locDelegateHasEmptyCode \\
						% 	\end{array} \right]
						% 	\vspace{2mm} \\
						% 	\locRecipientIsDelegatedAndDelegateIsDelegated & \define &
						% 	\left[ \begin{array}{clr}
						% 		\cdot & \locRecipientIsDelegated \\
						% 		\cdot & \locDelegateIsDelegated  \\
						% 	\end{array} \right]
						% 	\\
						% \end{array} \right.
					\]
					and define and derive some compound shorthands as follows
					\[
						\left\{ \begin{array}{lcl}
							\locRecipientHasEmptyCode    & \define & 1 - \locRecipientHasNonemptyCode                                                \\
							\locDelegateHasEmptyCode     & \define & 1 - \locDelegateHasNonemptyCode                                    \vspace{2mm} \\
							\locRecipientIsntDelegated   & \define & 1 - \locRecipientIsDelegated                                                    \\
							% \locDelegateIsntDelegated    & \define & 1 - \locDelegateIsDelegated                                        \vspace{2mm} \\
							\locRecipientHasNonemptyCode & \define & \accHasCode     _{i + \roffTxSkipUserRecipientAccount}                          \\
							\locDelegateHasNonemptyCode  & \define & \accHasCode     _{i + \roffTxSkipUserRecipientOrDelegationAccount} \vspace{2mm} \\
							\locRecipientIsDelegated     & \define & \accIsDelegated _{i + \roffTxSkipUserRecipientAccount}                          \\
							% \locDelegateIsDelegated      & \define & \accIsDelegated _{i + \roffTxSkipUserRecipientOrDelegationAccount}              \\
						\end{array} \right.
					\]
					\saNote{}
					% The terms marked with $\undefinedStar$ may be omitted from the implementation.
					There is no need to require accounts have nonempty code or have given code size to
					prevent invalid values of ``\accIsDelegated''.
					These are extraneous conditions given
					constraint~(\ref{hub: account: account delegation: newly delegated accounts have code and code size equal to 23}) and
					constraint~(\ref{hub: account: account delegation: delegated accounts have code and code size equal to 23}).
				\item
					\If   $\locIsDeploymentTransaction = 1$
					\Then $\txInitCodeSize _{i + \roffTxSkipUserTransaction} = 0$
			\end{enumerate}
		\item[\underline{Justifying $\txFinalRefundCounter$:}]
			we impose that $\txFinalRefundCounter _{i + \roffTxSkipUserTransaction} = \refund _{i}$

			\saNote{}
			Starting with \textsc{Prague} and more specifically with
			\cite{EIP-7702}, type-4 transactions which require no code execution
			still accrue refunds due to successful delegation operations.
		\item[\underline{Justifying $\txInitialBalance$:}]
			we impose that $\txInitialBalance _{i + \roffTxSkipUserTransaction} = \accBalance _{i + \roffTxSkipUserSenderAccount}$
		\item[\underline{Justifying \txStatusCode{}:}]
			we impose that $\txStatusCode _{i + \roffTxSkipUserTransaction} = 1$
		\item[\underline{Justifying \txNonce{}:}]
			we impose
			\[
				\accNonce                              _{i + \roffTxSkipUserSenderAccount} =
				\txNonce                               _{i + \roffTxSkipUserTransaction} +
				\txNumberOfSuccessfulSenderDelegations _{i + \roffTxSkipUserTransaction}
			\]
		\item[\underline{Justifying $\txLeftoverGas$:}]
			we impose that $\txLeftoverGas _{i + \roffTxSkipUserTransaction} = \txInitialGas _{i + \roffTxSkipUserTransaction}$
	\end{description}
	The following constraint is true up to the \textsc{Cancun} hardfork, but fails beyond that point.
	Specifically \cite{EIP-7623}, which is part of the \textsc{Prague} hardfork,
	imposes a floor price on call data
	which invalidates the upcoming ``sanity-check'' constraint.
	\begin{description}
		\item[\underline{Justifying $\txEffectiveRefund$:}]
			if the network runs on the \textsc{Cancun}-\evm{},
			we may impose the following sanity check:
			\[ \txEffectiveRefund _{i + \roffTxSkipUserTransaction} = \txLeftoverGas _{i + \roffTxSkipUserTransaction} \quad (\sanityCheck) \]
	\end{description}
	\iffalse
	I had to switch the order of phases from

	,   (TX_WARM) → (TX_AUTH) → (TX_SKIP ∨ TX_INIT) → ...
	,,,,,,,,,,,,,,,,,,,,,,,,,,,,,,,,,,,,,,,,,,,,,,,,,,,,,

	to 

	,   (TX_AUTH) → (TX_WARM) → (TX_SKIP ∨ TX_INIT) → ...
	,,,,,,,,,,,,,,,,,,,,,,,,,,,,,,,,,,,,,,,,,,,,,,,,,,,,,

	the reason for this switch is that it is easier to determine tx.requiresEvmExecution(), which decides whether we trace the TX_WARM phase, and whether we do TX_SKIP or TX_INIT, _after_ the TX_AUTH phase, given that the recipient account may have been delegated to a smart contract in the delegation list processing. Or the opposite.

	The real annoyance is that

	,   tx.requiresEvmExecution()
	,,,,,,,,,,,,,,,,,,,,,,,,,,,,,

	can no longer be determined from the transaction and the initial world state alone.
	\fi
	%
	% We further impose that every \textbf{set code transaction} i.e. every \textbf{type 4 transaction}
	% begin with the processing of its delegation list.
	% \begin{description}
	% 	\item[\underline{Set code transactions must incorporate a \txAuth{}-phase:}]
	% 		\If $\txTypeSupportsDelegationLists _{i} = \true$ \Then
	% 		\[
	% 			\left\{ \begin{array}{lcl}
	% 				\txAuth          _{i - 1} & = & \true \\
	% 				\peekTransaction _{i - 1} & = & \true \\
	% 			\end{array} \right.
	% 		\]
	% \end{description}
	% \saNote{}\label{hub: tx skip: user: transation processing: transactions supporting delegation lists start must contain a TX AUTH phase}
	% Compare with
	% constraint~(\ref{hub: initialization phase: common constraints: transactions supporting delegation lists must contain a TX AUTH phase}).
	% Also recall
	% constraint~(\ref{hub: authorization phase: perspectives: the transation type must support delegation lists})
	% and the associated discussion.

